% ============================================================
% SECTION 21 — THE POLICY ENGINE
% ============================================================

\section{The Policy Engine}

\begin{tcolorbox}[summarybox={\iconShield[1.5] Deterministic Business Logic}]
\color{textdark}
The Policy Engine (\filepath{app/policy/rules.py}) is the true decision-maker in RecoverAI V2. While the LLM provides a "recommended action," it is treated strictly as an untrusted suggestion. The Policy Engine evaluates that suggestion against hard, mathematical business rules and risk thresholds.

If the LLM recommends an action that violates a policy (either due to a hallucination, prompt injection, or simply a lack of business context), the Policy Engine mathematically \textbf{downgrades} the action to a safe fallback.
\end{tcolorbox}

\vspace{0.8cm}

% --- THE EVALUATION HIERARCHY ---
\subsection{The Evaluation Hierarchy}

The engine evaluates rules in a strict, descending hierarchy. A failure at any higher level immediately terminates the evaluation and forces a downgrade.

\begin{center}
\begin{tikzpicture}[
    node distance=0.8cm,
    rule/.style={rectangle, draw=primary, thick, fill=bglight, minimum width=6cm, minimum height=0.8cm, font=\small\bfseries, align=center},
    arrow/.style={-{Stealth[scale=1.1]}, thick, draw=textdark}
]

    \node[rectangle, draw=textdark, thick, dashed, minimum width=6cm, minimum height=0.8cm, font=\small] (llm) {LLM Recommendation: \code{RETRY\_PAYMENT}};

    \node[rule, below=0.8cm of llm] (r1) {1. Global Retry Budget Check};
    \node[rule, below=of r1] (r2) {2. Financial Amount Thresholds};
    \node[rule, below=of r2] (r3) {3. Absolute Risk Blocklist};
    
    \node[rectangle, draw=success, thick, fill=success!10, below=1cm of r3, minimum width=6cm, minimum height=0.8cm, font=\small\bfseries] (pass) {Final Action: \code{RETRY\_PAYMENT}};
    
    \node[rectangle, draw=danger, thick, fill=danger!10, right=1.5cm of r2, minimum width=3.5cm, minimum height=0.8cm, font=\small\bfseries, align=center] (fail) {Downgrade to\\ \code{STOP\_AUTOMATION}};

    \draw[arrow] (llm) -- (r1);
    \draw[arrow] (r1) -- (r2);
    \draw[arrow] (r2) -- (r3);
    \draw[arrow] (r3) -- (pass);
    
    \draw[arrow, draw=danger] (r1.east) -| (fail.north);
    \draw[arrow, draw=danger] (r2.east) -- (fail.west);
    \draw[arrow, draw=danger] (r3.east) -| (fail.south);

\end{tikzpicture}
\end{center}

\vspace{0.8cm}

% --- RULE DEFINITIONS ---
\subsection{Rule Definitions}

\subsubsection{Rule 1: Global Retry Budget}
The system enforces a hard limit on how many times a transaction can be retried (\code{MAX\_RETRIES}, default 2).
\begin{itemize}
    \item \textbf{Logic:} If \code{attempt\_count $\ge$ MAX\_RETRIES}, the engine forces \code{STOP\_AUTOMATION}.
    \item \textbf{Why:} Prevents infinite retry loops and runaway gateway fees. The LLM does not know how many times the system has already retried; the Policy Engine tracks this via the database.
\end{itemize}

\vspace{0.4cm}

\subsubsection{Rule 2: Financial Amount Thresholds}
The system enforces a maximum automatic retry limit (\code{MAX\_AUTO\_ACTION\_AMOUNT}, default \$50.00 / \code{5000} minor units).
\begin{itemize}
    \item \textbf{Logic:} If the transaction \code{amount > 5000} AND the LLM recommends \code{RETRY\_PAYMENT}, the engine downgrades the action to \code{CREATE\_ESCALATION} (flagging it for human review).
    \item \textbf{Why:} High-value transactions pose a higher risk of chargebacks or customer anger if retried aggressively.
\end{itemize}

\vspace{0.4cm}

\subsubsection{Rule 3: Absolute Risk Blocklist}
Certain gateway failure codes indicate that retrying is mathematically impossible or legally risky.
\begin{itemize}
    \item \textbf{Logic:} If \code{failure\_reason} is in \code{["fraud\_suspected", "stolen\_card", "account\_closed"]}, the engine forces \code{STOP\_AUTOMATION}.
    \item \textbf{Why:} If an attacker uses Prompt Injection to convince the LLM to retry a stolen card, this rule acts as the final deterministic blockade.
\end{itemize}

\vspace{0.8cm}

% --- THE ENGINE IMPLEMENTATION ---
\subsection{Engine Implementation}

\begin{tcolorbox}[codebox={Policy Evaluation (\filepath{app/policy/rules.py})}]
\begin{lstlisting}[language=Python]
def evaluate_action(
    transaction: Transaction,
    llm_recommendation: str,
    attempt_count: int
) -> str:
    # 1. Start with the LLM's suggestion
    action = llm_recommendation
    
    # 2. Rule 1: Budget
    if attempt_count >= settings.MAX_RETRIES:
        return "STOP_AUTOMATION"
        
    # 3. Rule 2: Amount Threshold
    if action == "RETRY_PAYMENT" and transaction.amount > settings.MAX_AUTO_ACTION_AMOUNT:
        return "CREATE_ESCALATION"
        
    # 4. Rule 3: Risk Blocklist
    fatal_codes = ["fraud", "stolen_card", "lost_card"]
    if transaction.failure_reason in fatal_codes:
        return "STOP_AUTOMATION"
        
    # 5. Passed all rules
    return action
\end{lstlisting}
\end{tcolorbox}

\begin{tcolorbox}[successbox={Deterministic Safety}]
By isolating the non-deterministic AI (LLM) behind a wall of deterministic Python (\code{if/else} statements), the system guarantees that business rules are always obeyed, satisfying compliance and security requirements.
\end{tcolorbox}
